\documentclass[tikz,border=3pt]{standalone}
\input{tikz-tcee}
\begin{document}
\hyphenpenalty=10000 \exhyphenpenalty=10000
% Triángulo de Gordon: de dónde viene la inflación
\begin{tikzpicture}[tcee, font=\footnotesize,
  v/.style={draw=#1, fill=#1!14, rounded corners=2pt, line width=0.7pt, text width=3.4cm, align=center, inner sep=3pt}]
  \fill[dicinflacion!6] (0,2.15) -- (-2.75,-0.6) -- (2.75,-0.6) -- cycle;
  \draw[thick, dicreferencia] (0,2.15) -- (-2.75,-0.6) -- (2.75,-0.6) -- cycle;
  \node[font=\large\bfseries] at (0,0.35) {$\pi$};
  \node[v=dicfbcf] at (0,2.55) {\textbf{Demanda}\\brecha $+$, crédito y M3 $\uparrow$, estímulo fiscal};
  \node[v=dicenergia] at (-2.65,-1.15) {\textbf{Costes} (oferta)\\CLU, energía, depreciación, márgenes};
  \node[v=dicservicios] at (2.65,-1.15) {\textbf{Inercia}\\expectativas, indexación, segunda ronda};
\end{tikzpicture}
\end{document}
